% !TEX root = ../main.tex

\chapter{Calculating derived intersections}
\label[chapter]{chap:Derived_Zero_Loci_And_Tangent_Complexes}

A smooth map $f\colon U\to\R^r$ determines a derived manifold
\[
Z^{\mathrm{der}}(f)=U\times_{\R^r}\{0\}.
\]
The preceding chapter guarantees that this pullback exists. We now ask
what its functions remember about the equations defining it. To calculate
them, we need an algebraic presentation of the pullback.
The resulting algebra has one generator for each equation, together with a
differential recording that equation. Products of the generators keep track
of relations among the equations.

Write $f=(f_1,\ldots,f_r)$. Under the
contravariant description of derived manifolds by animated $\Cinfty$-rings,
the defining pullback becomes the pushout
\[
\begin{tikzcd}
 \Cinfty(\R^r) \rar["f^*"] \dar["0^*"'] \drar[pushout]
   & \Cinfty(U) \dar\\
 \R \rar & \mathcal O(Z^{\mathrm{der}}(f)).
\end{tikzcd}
\]
Thus
\begin{equation}\label{eq:Derived_Zero_Locus_Animated_Pushout}
 \mathcal O(Z^{\mathrm{der}}(f))
 \cong \Cinfty(U)\mathbin{\otimes^{\mathbb L}_{\infty,\Cinfty(\R^r)}}\R.
\end{equation}
Here the tensor notation denotes a pushout of animated smooth rings.
The corresponding ordinary pushout is
$\Cinfty(U)/(f_1,\ldots,f_r)$. Computing that quotient first would discard
the higher homotopy of the animated pushout.

There are already two different classical objects in this discussion.
The ring $\Cinfty(U)/(f_1,\ldots,f_r)$ records the equations, whereas the
set $f^{-1}(0)$ only records their common real solutions. For example,
$x$ and $x^2$ have the same zero set but give different quotients.
The derived construction adds a third level: homotopy groups beyond
$\pi_0$.

We first introduce smooth differential graded algebras as a tool for
computing animated smooth rings, then construct the Koszul model of a
zero locus. Comparing a transverse equation, a multiple zero, and a
repeated equation will make these three levels visible. The same
construction will then apply to families and to sections of vector bundles.

\section{Smooth differential graded algebras}
\label[section]{sec:Smooth_Dold_Kan_Bridge}

Smooth differential graded algebras make the higher homotopy of an
animated smooth ring accessible through homology. We use homological
grading, as in Steffens: the differential lowers degree by one. Our
connective dg-algebra models are concentrated in nonnegative degrees.
After describing their smooth structure, we state the comparison theorem
that justifies using them to compute derived pushouts.

\begin{definition}[Smooth dg-algebra]
\label[definition]{def:Smooth_Dg_Algebra}
A \emph{connective $\Cinfty$-dg-algebra} is a commutative differential
graded $\R$-algebra
\[
 \cdots\longrightarrow A_2\xrightarrow{d}A_1
 \xrightarrow{d}A_0
\]
together with a $\Cinfty$-ring structure on $A_0$ whose underlying
$\R$-algebra structure is the given one. A morphism preserves the differential,
the multiplication, and the smooth operations on degree zero.
\end{definition}

For homogeneous elements, commutativity and the differential rule mean
\[
 ab=(-1)^{|a||b|}ba,\qquad
 d(ab)=d(a)b+(-1)^{|a|}a\,d(b).
\]
In particular, an odd-degree generator has square zero over $\R$.
The smooth operations are supplied in degree zero. There is no operation
$f(a_1,\ldots,a_n)$ for arbitrary nonhomogeneous elements in this definition.

This is the convention of the foundations paper \cite[Definition~3.3.1]
{Steffens_Derived_Cinfty_Geometry_I}. Because $d$ vanishes on $A_0$, there
is no additional chain-rule condition to impose on its degree-zero elements.

\begin{example}
An ordinary $\Cinfty$-ring is a smooth dg-algebra concentrated in degree
zero. More generally, for an ordinary smooth ring $A$ and elements
$f_1,\ldots,f_r\in A$, form
\[
 A[\varepsilon_1,\ldots,\varepsilon_r],
 \qquad |\varepsilon_i|=1,\qquad d\varepsilon_i=f_i.
\]
The notation denotes a free graded-commutative extension: as a graded
$A$-module it is $A\otimes_\R\Lambda^\bullet(\R^r)$, with
$\Lambda^q$ placed in degree $q$. Its differential is determined by the
displayed equations and the Leibniz rule. In particular,
\[
 d(\varepsilon_i\varepsilon_j)
 =f_i\varepsilon_j-f_j\varepsilon_i.
\]
The equality $d^2=0$ follows on generators and hence on the whole algebra.
\end{example}

The image of $d\colon A_1\to A_0$ is an ideal: for $a\in A_0$ and
$b\in A_1$, $a\,d(b)=d(ab)$. Therefore $H_0(A)$ is an ordinary
$\Cinfty$-ring. To see that all smooth operations descend, apply Hadamard's
lemma to a smooth function $\varphi\colon\R^n\to\R$:
\[
 \varphi(x)-\varphi(y)=\sum_i(x_i-y_i)\psi_i(x,y).
\]
Replacing each $x_i$ by an element congruent to $y_i$ modulo the ideal
changes the output by an element of the same ideal. The groups $H_i(A)$
are modules over $H_0(A)$, and their products form a graded algebra.

The passage from dg-algebras to animated smooth rings identifies maps
that induce isomorphisms on homology.

\begin{definition}
\label[definition]{def:Smooth_Dga_Quasi_Isomorphism}
A morphism of smooth dg-algebras is a \emph{quasi-isomorphism} if it induces
an isomorphism on every homology group.
\end{definition}

\begin{theorem}[Smooth Dold--Kan correspondence]
\label[theorem]{fact:Smooth_Dold_Kan_Comparison}
The localization of connective $\Cinfty$-dg-algebras at quasi-isomorphisms,
taken as an $\infty$-categorical localization, is equivalent to the
$\infty$-category of animated $\Cinfty$-rings. If $A_{\mathrm{an}}$
corresponds to a dg-algebra $A$, then
\[
 \pi_i(A_{\mathrm{an}})\cong H_i(A),\qquad i\geq0.
\]
\end{theorem}

We use this as a comparison theorem
\cite[Theorem~3.3.12]{Steffens_Derived_Cinfty_Geometry_I}.
It includes multiplicative and smooth structure. The ordinary Dold--Kan
correspondence for vector spaces alone does not supply that structure.

One way to describe the animated object attached to $A$ is through its
values on free smooth rings. In the localized category $D$ of smooth
dg-algebras, set
\[
 A_{\mathrm{an}}(\R^n)=\Hom_D(\Cinfty(\R^n),A).
\]
The functor is covariant in Euclidean spaces because smooth functions are
contravariant. Free smooth rings turn finite products of Euclidean spaces
into coproducts, so this functor preserves finite products. It is therefore
an animated $\Cinfty$-ring. The mapping objects in this formula are animae;
using only strict dg-algebra maps would lose the higher homotopies.

The theorem allows computations using dg-algebras provided we compute
their \emph{derived} pushouts. An arbitrary strict pushout of dg-algebras
need not have the required universal property. In the zero-locus
calculation, we will replace evaluation at the origin by a free resolution
before taking the strict pushout.

\section{The Koszul model of a zero locus}
\label[section]{sec:Koszul_Presentation_Derived_Zero_Locus}

We compute the pushout in
\Cref{eq:Derived_Zero_Locus_Animated_Pushout} in two stages. First we
resolve evaluation at the origin of $\R^r$. We then substitute the
functions defining our zero locus into this resolution.

Start with one coordinate $t$ and put $B=\Cinfty(\R)$. Consider the map
\[
 B[\varepsilon]\longrightarrow\R,\qquad
 |\varepsilon|=1,\quad d\varepsilon=t,
\]
sending $t$ and $\varepsilon$ to zero. Its underlying complex is
\[
 B\xrightarrow{\ t\ }B
\]
in degrees $1,0$. Multiplication by $t$ is injective: a smooth function
annihilated by $t$ vanishes away from the origin, hence everywhere.
Its cokernel is $\R$, since a smooth function vanishing at zero is divisible
by $t$. Thus this map is a quasi-isomorphism.

For $B=\Cinfty(\R^r)$, the same construction gives
\[
 K_B(t_1,\ldots,t_r)
 =\bigl(B[\varepsilon_1,\ldots,\varepsilon_r],
          d\varepsilon_i=t_i\bigr)\longrightarrow\R.
\]
The coordinate functions form a regular sequence. Indeed, quotienting by
the first $j$ coordinates gives the smooth functions on the remaining
coordinate subspace, and the next coordinate is a non-zero-divisor there.
Successively applying the one-variable calculation proves that the
augmentation is a quasi-isomorphism.

The map $B\to K_B(t_1,\ldots,t_r)$ is built by adjoining generators in
positive degree with prescribed differential. Together with the
augmentation to $\R$, it is a relative free resolution of evaluation.
Taking its strict pushout computes derived base change. The criterion
used here is \citet[Corollary~3.3.15]{Steffens_Derived_Cinfty_Geometry_I}:
this map is an isomorphism in
degree zero, is surjective on $H_0$, and is a relative cell extension
of underlying commutative dg-algebras. The resulting zero-locus
calculation is given in
\citet[Example~3.3.16]
{Steffens_Derived_Cinfty_Geometry_I}.

Now base change along $f^*\colon B\to A=\Cinfty(U)$. Each $t_i$ becomes
$f_i$, so the pushout is
\begin{equation}\label{eq:Koszul_Algebra_Function}
 K_A(f)=\bigl(A[\varepsilon_1,\ldots,\varepsilon_r],
               d\varepsilon_i=f_i\bigr).
\end{equation}

\begin{theorem}[Koszul model]
\label[theorem]{thm:Koszul_Presentation_Derived_Zero_Locus}
The smooth dg-algebra $K_A(f)$ presents the animated function ring of
$Z^{\mathrm{der}}(f)$. Consequently,
\[
 \pi_i\mathcal O(Z^{\mathrm{der}}(f))\cong H_i(K_A(f)).
\]
\end{theorem}

\begin{proof}
Replace the evaluation map $B\to\R$ in
\Cref{eq:Derived_Zero_Locus_Animated_Pushout} by the relative free
resolution $B\to K_B(t)$. Its strict pushout along $B\to A$ computes
the derived pushout and is exactly \Cref{eq:Koszul_Algebra_Function}.
Apply the smooth Dold--Kan correspondence.
\end{proof}

The construction resolves the coordinates on $\R^r$, which are a regular
sequence, before substituting the possibly singular functions $f_i$.
No regular-sequence assumption on the $f_i$ is needed. The
positive-degree homology can retain additional relations among them.

In particular,
\[
 H_0(K_A(f))=A/(f_1,\ldots,f_r).
\]
In degree $1$, a cycle $\sum_i a_i\varepsilon_i$ is a relation
$\sum_i a_if_i=0$. Boundaries include the relations
$f_i\varepsilon_j-f_j\varepsilon_i$. Thus higher homology measures
relations that survive after these elementary ones have been identified.
The algebra also retains the multiplication between such classes.

\section{Equations with the same zero set}

The preceding chapters computed ordinary quotient rings. We now return
to those equations to ask what survives in the positive degrees of the
Koszul model. In particular, a multiple zero need not have higher
homotopy of functions, whereas an excess equation can produce it.

\begin{example}[A transverse equation]
For $f(x)=x$, the model is
\[
 \Cinfty(\R)[\varepsilon],\qquad d\varepsilon=x.
\]
The resolution calculation shows that it is quasi-isomorphic to $\R$.
Thus the derived zero locus is the ordinary point.

More generally, if $f\colon U\to\R^r$ is transverse to zero, its components
can be completed locally to coordinates along the zero set. The same
Koszul calculation shows that the derived zero locus is the ordinary
manifold $f^{-1}(0)$. Away from the zero set, one of the functions is
invertible, and the Koszul complex is contractible. These local statements
recover the transverse-pullback property of derived manifolds.
\end{example}

\begin{example}[A multiple zero]
For $f(x)=x^2$, multiplication by $x^2$ is still injective on $\Cinfty(\R)$.
Hence the Koszul algebra has no positive homology, but
\[
 H_0(K(x^2))=\Cinfty(\R)/(x^2)\cong\R[\delta]/(\delta^2).
\]
The last identification follows from Taylor's formula:
\[
 g(x)=g(0)+xg'(0)+x^2r(x),
\]
with $r$ smooth. The smooth operations on the quotient are given by the
first-order Taylor formula. For instance,
$\varphi(a+b\delta)=\varphi(a)+\varphi'(a)b\delta$.

Thus $Z^{\mathrm{der}}(x^2)$ is not the ordinary point, even though its
animated ring is discrete. Nilpotents in $\pi_0$ can already distinguish
it. The same argument gives
$\Cinfty(\R)/(x^m)\cong\R[\delta]/(\delta^m)$ for $m\geq1$.
\end{example}

\begin{example}[A zero equation]
For the zero function on $\R$, the differential vanishes:
\[
 K(0)=\Cinfty(\R)[\varepsilon],\qquad d\varepsilon=0.
\]
Its degree-zero homology is $\Cinfty(\R)$, and its degree $1$
homology is another copy of $\Cinfty(\R)$. The classical truncation is
the line, while the derived object remembers that an equation was imposed
which vanishes identically.

Similarly, the self-intersection of $\{0\}\subset\R$ has function algebra
$\R[\varepsilon]$ with $d\varepsilon=0$. Its underlying topological space
is a point and its classical function ring is $\R$, but $\pi_1=\R$.
This is a local instance of the excess-intersection phenomena in
\citet[Example~2.7]{Spivak_Derived_Smooth_Manifolds}.
\end{example}

\begin{example}[A repeated equation]
The maps $x\mapsto x$ and $x\mapsto(x,x)$ have the same zero set and
the same ordinary quotient ring. For the second map the Koszul algebra is
\[
 \Cinfty(\R)[\varepsilon_1,\varepsilon_2],
 \qquad d\varepsilon_1=d\varepsilon_2=x.
\]
Set $\eta=\varepsilon_2-\varepsilon_1$. Then $d\eta=0$, and the algebra
becomes $K(x)\otimes_\R\Lambda(\eta)$. It is quasi-isomorphic to
$\R[\eta]$ with zero differential. The repeated equation produces a
nonzero $\pi_1$.
\end{example}

These examples separate three comparisons. A point and a multiple point
differ in their classical function rings. An ordinary point and a derived
self-intersection differ in higher homotopy. Different lists generating
the same ordinary ideal can define different derived zero loci, because
the derived construction remembers the chosen map to the space of equations.

\section{Families of equations}

The model also lets us follow an equation as its parameters vary.
Even when the total zero locus is a manifold, its fibres can acquire
multiplicity or higher homotopy.

Let $S$ be a parameter manifold and let
$f\colon S\times U\to\R^r$ be smooth. The relative equation defines a
derived object over $S$:
\[
 X=(S\times U)\times_{S\times\R^r}(S\times\{0\}).
\]
Its Koszul model has coefficients in $\Cinfty(S\times U)$ and
differential $d\varepsilon_i=f_i(s,x)$.

For a smooth map $g\colon T\to S$, pullback gives the same construction
with $f(g(t),x)$. This follows by pasting the defining pullback squares.
In particular, the fibre at $s\in S$ is the derived zero locus of $f_s$.
The fibre is a derived pullback even when the total object is an ordinary
manifold.

For example, consider $f(s,x)=x^2-s$. The total zero locus is the ordinary
parabola $s=x^2$, since the derivative in the $s$-direction is $-1$.
Over $s>0$ its fibre consists of two points; over $s<0$ it is empty.
The derived fibre at $s=0$ has function ring $\R[\delta]/(\delta^2)$.
The parameter map from the parabola to $S$ has a critical point there.

A second family, $f(s,x)=sx$, illustrates the appearance of higher
homotopy in a fibre. For $s\neq0$, the fibre is an ordinary point.
For $s=0$, its Koszul model is
$\Cinfty(\R)[\varepsilon]$ with zero differential. Although multiplication
by $sx$ on $\Cinfty(\R^2)$ is injective, substitution $s=0$ destroys
injectivity. Derived base change retains this change in the equation.

\section{Zero loci of vector bundle sections}
\label[section]{sec:Supplement_Vector_Bundle_Koszul_Formula}

The same model works without a chosen list of global equations. Let
$p\colon E\to U$ be a vector bundle of rank $r$, and let $s$ be a section.
Its derived zero locus is the pullback of $s$ and the zero section as maps
from $U$ to $E$. Define a sheaf of smooth dg-algebras on $U$ by
\[
 \mathcal K(s)_q
   =\Gamma(-,\Lambda^qE^\vee),\qquad d=\iota_s.
\]
Here $\iota_s$ contracts an alternating form with $s$. In a local frame,
this is the Koszul algebra already calculated. Changes of frame respect
contraction, so the local algebras glue.

\begin{proposition}
\label[proposition]{prop:Vector_Bundle_Koszul_Presentation}
The sheaf $\mathcal K(s)$ presents the functions of the derived zero locus
of $s$, supported on its ordinary zero set. In particular, for the zero
section its differential vanishes and its homotopy sheaves are
$\Lambda^qE^\vee$ in degree $q$.
\end{proposition}

\begin{proof}
Trivialize $E$ on an open cover and apply the Koszul model theorem.
On the complement of the zero set, contraction has a local contracting
homotopy: choose $\lambda\in E^\vee$ with $\lambda(s)=1$ and use exterior
multiplication by $\lambda$. The frame-independent contraction formula
identifies the models on overlaps, giving the stated sheaf.
\end{proof}

This is the bundle form of the calculation in
\cite[Example~3.3.17]{Steffens_Derived_Cinfty_Geometry_I}.
It explains how the local excess directions of a self-intersection fit
into a vector bundle rather than a fixed vector space.

\section{Further reading}

The following observations explain the role of regular sequences in the
examples and the effect of allowing generators in higher degrees.
The exercises give further comparisons between presentations of zero loci.

A regular sequence $f_1,\ldots,f_r$ in an ordinary
smooth ring has Koszul homology concentrated in degree zero. One proves
this inductively: adjoining the last generator forms the cone of
multiplication by $f_r$ on the preceding Koszul complex. If the preceding
homology is $A/(f_1,\ldots,f_{r-1})$ and $f_r$ is a non-zero-divisor
there, the cone has only the desired quotient as its degree-zero
homology. This explains both the transverse case and the multiple-zero
example. Concentration in degree zero is weaker than smoothness of the
resulting ordinary ring.

The smooth dg-algebra model also permits generators in degrees above $1$.
For example, a free generator $z$ in degree $2$ with $dz=0$ gives
$\R[z]$. Since $z$ has even degree, its powers need not vanish:
$H_{2j}(\R[z])=\R z^j$ for every $j\geq0$.
A finite list of dg-generators therefore does not imply that the
function algebra has only finitely many nonzero homotopy groups.
The boundedness of a tangent complex is a separate issue, treated next.

\begin{exercise}
Compute the Koszul homology of $(x,0)\colon\R\to\R^2$ and compare it
with the repeated-equation example.
\end{exercise}

\begin{exercise}
For $f=(x_1,\ldots,x_q,0,\ldots,0)\colon\R^n\to\R^r$, show that the
Koszul algebra is quasi-isomorphic to
$\Cinfty(\R^{n-q})[\eta_1,\ldots,\eta_{r-q}]$ with zero differential.
Interpret the surviving generators as excess equations.
\end{exercise}

\begin{exercise}
Show directly that adjoining a new coordinate $y$ and the equation $y=0$
does not change the derived zero locus. Compare this with adjoining the
identically zero equation.
\end{exercise}
